% -----------------------------------------------------------
\section{1837}
% -----------------------------------------------------------
	\label{EMS-1837}
	
	% ----------------------
	\subsection{Joseph Smith}
	% ----------------------
		\label{EMS-1837-JS}

		% ----------------------
		\subsubsection{Introduction to 1837}
		% ----------------------
	
			sdf
		
		% -----
		\subsubsection{Discourse, 6 April 1837}
		% -----
			\label{EMS-1837-JS-6-APR-1837}
			% \label{EMS-1830-JS-DAY-3LETTERMONTH-YEAR}
			
			\begin{description}
					\item[Print Sources]
				\end{description}

					\begin{tabular}{p{0.5in}p{1in}p{0.5in}p{1in}}
					JSP	 	& D5:352--357		& JSR 		& --- \\
					DC	 	& ---				& HC 		& 2:477--479 \\
					HDDC 	& ---				& TCHC 		& 2:461--463 \\
					TPJS 	& 111--113			& 	 		&  \\
					\end{tabular}
					% HDDC = woodford's DC diss; JSR = Marquardt's JS Revelations; TCHC = Vogel HC
				
				\begin{description}
					\item[Online Access] \href{https://www.josephsmithpapers.org/paper-summary/discourse-6-april-1837/1}{Here}
					% \item[Analysis] \hyperref[XX]{Here}
				\end{description}
				
				\begin{description}
					\item[Background]
				\end{description}
				
					% It might also be useful to give a two sentence context of the period: e.g., "This speech was given in Nauvoo to a small congregation one month prior to the destruction of the Nauvoo Expositor. These and other tensions may have been on the prophet's mind when he gave the speech."
				
				\begin{description}
					\item[Original Source]
				\end{description}
				
					% see also messenger and advocate 3:7 pages 486--489
				
					% brief description of original source material, with reference to JSP or somewhere else for more info
					% Background material: it might be helpful to have a note that says whether a given document was presented as a speech, was written in a journal, etc.
					% Also a comment here about how reliable the source might be, or what shape it is in, if there are large omissions or parts that are hard to understand, and so on.
				
				\begin{description}
					\item[Text]
				\end{description}
				
						\hypertarget{EMS-1837-JS-6-APR-1837-a}%
					Joseph
						\marginnote{a}%
					Smith jr. rose and spoke on the subject of the Priesthood. The Melchisedec High priesthood, he said was no other then the priesthood of the Son of God. There are certain ordinances which belong to the priesthood, and certain results flow from it.
					
					The presidents, or presidency are over the church, and revelations of the mind and will of God to the church are to come through the presidency. This is the order of heaven and the power and privilege of this priesthood. It is also the privilege of any officer in this church, to obtain revelations so far as relates to his particular calling or duty in the church. All are bound by the principles of virtue and happiness, but one great privilege of this priesthood is to obtain revelations, as before observed, of the mind and will of God. It is also the privilege of the Melchisedec priesthood, to reprove, rebuke and admonish, as well as to receive revelations.
					
						\hypertarget{EMS-1837-JS-6-APR-1837-b}%
					He
						\marginnote{b}%
					here remarked something concerning the will of God, and said, that what God commanded, the one half of the church would condemn.--- A high Priest, is a member of the same Melchisedec priesthood, with the presidency, but not of the same power or authority in the church. The seventies are also members of the same priesthood, are a sort of travelling council, or priesthood, and may preside over a church or churches until a high priest can be had. The seventies are to be taken from the quorum of elders and are not to be high priests. They are subject to the direction and dictation of the twelve, who have the keys of the ministry. All are to preach the gospel, by the power and influence of the Holy Ghost, and no man, said he, can preach the gospel without the Holy Ghost.
					
						\hypertarget{EMS-1837-JS-6-APR-1837-c}%
					The
						\marginnote{c}%
					Bishop was a high priest, and necessarily so, because he is to preside over that particular branch of the church affairs that are denominated the lesser priesthood, and because we have no direct lineal descendent of Aaron to whom it would of right belong. He remarked that this was the same, or a branch of the same priesthood; and illustrated his position by the figure of the human body, which has dfferent members, which have different offices to perform: all are necessary in their place, and the body is not complete without all the members. From a view of the requirements of the servants of God to preach the gospel, he remarked that few were qualified even to be priests, and if a priest understood his duty, his calling and ministry and preached by the Holy Ghost, his enjoyment is as great as if he were one of the presidency; and his services are necessary in the body, as are also those of teachers and deacons. Therefore in viewing the church as whole, we may strictly denominate it one priesthood.
					
						\hypertarget{EMS-1837-JS-6-APR-1837-d}%
					He
						\marginnote{d}%
					remarked that he rebuked and admonished his brethren frequently, and that because he loved them; not because he wished to incur their displeasure or mar their happiness.
					
					Such a course of conduct was not calculated to gain the good will of all, but rather the ill will of many, and thereby the situation in which he stood was an important one. So you see, brethren the higher the authority, the greater the difficulty of the station. But these rebukes and admonitions became nccssary from the perverseness of brethren, for their temporal as well as spiritual welfare. They actually constituted a part of the duties of his station and calling.
					
						\hypertarget{EMS-1837-JS-6-APR-1837-e}%
					Others
						\marginnote{e}%
					had other duties to perform that were important and far less enviable, and might be just as good, like the feet or hands in their relation to the human body, neither could claim priority, or say to the other I have no need of you. After all that has been said the greatest duty and the most important is, to preach the gospel.
					
					He then alluded to the temporal affairs of the church in this place, stating the causes of the embarrassments of a pecuniary nature that were now pressing upon the heads of the church. He observed they began poor, were needy, destitute, and were truly afflicted by their enemies; yet the Lord commanded them to go forth and preach the gospel, to sacirfice their time, their talents, their good name and jeopardize their lives, and in addition to this, they were to build a house for the Lord, and prepare for the gathering of the saints.
					
						\hypertarget{EMS-1837-JS-6-APR-1837-f}%
					Thus
						\marginnote{f}%
					it was easy to see this must involve them. They had no temporal means in the beginning commensurate with such an undertaking, but this work must be done, this place had to be built up. He further remarked that it must yet be built up, that more houses must be built. He observed that large contracts had been entered into for land on all sides where our enemies had signed away their right. We are indebted to them to be sure, but our brethren abroad have only to come with their money, take these eontracts [contracts], relieve their brethren of the pecuniary embarrassments under which they now labor, and procure for themselves a peaceable place of rest among us. He then closed at about 4 P. M. by uttering a prophesy saying this place must be built up, and would be built up, and that every brother that would take hold and help secure and discharge those contracts that had been made, should be rich.
					
		% -----
		\subsubsection{Revelation, 23 July 1837 [D\&C 112]}
		% -----
			\label{EMS-1837-JS-23-JUL-1837}
			% \label{EMS-1830-JS-DAY-3LETTERMONTH-YEAR}
			
			\begin{description}
					\item[Print Sources]
				\end{description}

					\begin{tabular}{p{0.5in}p{1in}p{0.5in}p{1in}}
					JSP	 	& D5:410--417; J1:306--309		& JSR 		& 281--283 \\
					DC	 	& Section 112		& HC 		& 2:499--501 \\
					HDDC 	& 1475--1494		& TCHC 		& 2:478--479 \\
					\end{tabular}
					% HDDC = woodford's DC diss; JSR = Marquardt's JS Revelations; TCHC = Vogel HC
				
				\begin{description}
					\item[Online Access] \href{https://www.josephsmithpapers.org/paper-summary/revelation-23-july-1837-dc-112/1}{Here}
					% \item[Analysis] \hyperref[DC112]{Here}
				\end{description}
				
				\begin{description}
					\item[Background]
				\end{description}
				
					% It might also be useful to give a two sentence context of the period: e.g., "This speech was given in Nauvoo to a small congregation one month prior to the destruction of the Nauvoo Expositor. These and other tensions may have been on the prophet's mind when he gave the speech."
				
				\begin{description}
					\item[Original Source]
				\end{description}
				
					% brief description of original source material, with reference to JSP or somewhere else for more info
					% Background material: it might be helpful to have a note that says whether a given document was presented as a speech, was written in a journal, etc.
					% Also a comment here about how reliable the source might be, or what shape it is in, if there are large omissions or parts that are hard to understand, and so on.
				
				\begin{description}
					\item[Text]
				\end{description}
				
					\begin{tightcenter}
					A Revelation given Kirtland July 23\supers{rd.} 1837.
					\end{tightcenter}

					The word of the Lord unto Thomas, B. Marsh concerning the twelve Apostles of the Lamb.

					Verily thus saith the Lord unto you my servant Thomas, I have heard thy prayers and thine alms have come up as a memorial before me in behalf of those thy brethren who were chosen to bear testimony of my name and to send it abroad among all nations, kindreds, tongues and people and ordained through the instrumentality of my servants.

					Verily I say unto you there have been some few things in thine heart and with thee, with which I the Lord was not well pleased; nevertheless inasmuch as thou hast abased thyself thou shalt be exalted: therefore all thy sins are forgiven thee. Let thy heart be of good cheer before my face, and thou shalt bear record of my name, not only unto the Gentiles, but also unto the Jews; and thou shalt send forth my word unto the ends of the earth.

					Contend thou therefore morning by morning, and day after day let thy warning voice go forth; and when the night cometh let not the inhabitants of the earth slumber because of thy speech. Let thy habitation be known in Zion, and remove not thy house, for I the Lord have a great work for \sout{you} <thee> to do, in publishing my name among the children of men, therefore gird up your loins for the work. Let your feet be shod also for thou art chosen, and thy path lyeth among the mountains and among many nations, and by thy word many high ones shall be brought low; and by thy word many low ones shall be exalted, thy voice shall be a rebuke unto the transgressor, and at thy rebuke let the tongue of the slanderer cease its perverseness. Be thou humble and the Lord thy God shall lead thee by the hand and give thee an answer to thy prayers, I know thy heart and have heard thy prayers concerning thy brethren. Be not partial towards them in love above many others, but let your love be for them as for yourself, and let your love abound unto all men and unto all who love my name. And pray for your brethren of the twelve. Admonish them sharply for my name's sake, and let them be admonished for all their sins, and be ye faithful before me unto my name; and after their temptations and much tribulation behold I the Lord will feel after them, and if they harden not their hearts and stiffen not their necks against me they shall be converted and I will heal them.

					Now I say unto you, and what I say unto you, I say unto all the twelve. Arise and gird up your loins, take up your cross, follow me, and feed my sheep. Exalt not yourselves; rebel not against my servant Joseph for Verily I say unto you I am with him and my hand shall be over him; and the keys which I have given him, and also to youward shall not be taken from him untill I come.

					Verily I say unto you my servant Thomas, thou art the man whom I have chosen to hold the keys of my kingdom (as pertaining to the twelve) abroad among all nations, that thou mayest be \sout{thy} my servant to unlock the door of the kingdom in all places where my servant Joseph, and my servant Sidney [Rigdon], and my servant Hyrum [Smith], cannot come, for on them have I laid the burden of all the Churches for a little season: wherefore whithersoever they shall send you, go ye, and I will be with you and in whatsoever place ye shall proclaim my name an effectual door shall be opened unto you that they may receive my word. Whosoever receiveth my word receiveth me, and whosoever receiveth me receiveth those (the first presidency) whom I have sent, whom I have made counsellors for my name's sake unto you. And again I say unto you, that whosoever ye shall send in my name, by the voice of your brethren the twelve, duly recommended and authorized by you, shall have power to open the door of my kingdom unto any nation whithersoever ye shall send them, inasmuch as they shall humble themselves before me and abide in my word, and hearken to the voice of my spirit.

					Verily verily! I say unto you, darkness covereth the earth and gross darkness the \sout{people} minds of the people, and all flesh has become corrupt before my face! Behold vengeance cometh speedily upon the inhabitants of the earth. A day of wrath! A day of burning! A day of desolation! Of weeping! Of mourning and of lamentation! And as a whirlwind it shall come upon all the face of the earth saith the Lord. And upon my house shall it begin and from my house shall it go forth saith the Lord. First among those among you saith the Lord; who have professed to know my name and have not known me and have blasphemed against me in the midst of my house saith the Lord

					Therefore see to it that you trouble not yourselves concerning the affairs of my Church in this place saith the Lord, but purify your hearts before me, and then go ye into all the world and preach my gospel unto every creature who have not received it and he that believeth and is baptized shall be saved, and he that believeth not, and is not baptized shall be damned. For unto you (the twelve) and those (the first presidency) who are appointed with you to be your counsellors and your leaders, is the power of this priesthood given for the last days and for the last time, in the which is the dispensation of the fulness of times: which power you hold in connection with all those who have received a dispensation at any time from the beginning of the creation, for verily I say unto you the keys of the dispensation which ye have received have came down from the fathers; and last of all being sent down from heaven unto you. Verily I say unto you, Behold how great is your calling.

					Cleanse your hearts and your garments, lest the blood of this generation be required at your hands. Be faithful untill I come for I come quickly and my reward is with me to recompense every man according as his work shall be! I am Alpha and Omega. Amen.
					
		% -----
		\subsubsection{Minutes, 3 September 1837}
		% -----
			\label{EMS-1837-JS-3-SEP-1837}
			% \label{EMS-1830-JS-DAY-3LETTERMONTH-YEAR}
			
			\begin{description}
					\item[Print Sources]
				\end{description}

					\begin{tabular}{p{0.5in}p{1in}p{0.5in}p{1in}}
					JSP	 	& D5:420--425		& JSR 		& --- \\
					DC	 	& ---				& HC 		& 2:509--510 \\
					HDDC 	& ---				& TCHC 		& 2:485--487 \\
					\end{tabular}
					% HDDC = woodford's DC diss; JSR = Marquardt's JS Revelations; TCHC = Vogel HC
				
				\begin{description}
					\item[Online Access] \href{https://www.josephsmithpapers.org/paper-summary/minutes-3-september-1837/1}{Here}
					% \item[Analysis] \hyperref[XX]{Here}
				\end{description}
				
				\begin{description}
					\item[Background]
				\end{description}
				
					% It might also be useful to give a two sentence context of the period: e.g., "This speech was given in Nauvoo to a small congregation one month prior to the destruction of the Nauvoo Expositor. These and other tensions may have been on the prophet's mind when he gave the speech."
				
				\begin{description}
					\item[Original Source]
				\end{description}
				
					% brief description of original source material, with reference to JSP or somewhere else for more info
					% Background material: it might be helpful to have a note that says whether a given document was presented as a speech, was written in a journal, etc.
					% Also a comment here about how reliable the source might be, or what shape it is in, if there are large omissions or parts that are hard to understand, and so on.
				
				\begin{description}
					\item[Text]
				\end{description}
				
						\hypertarget{EMS-1837-JS-3-SEP-1837-a}%
					...President
						\marginnote{a}%
					Rigdon then arose, \& made an address of conciderable length, showing the starting point or cause of all the difficulty of Elders Boyngton \& Johnson, he allso cautioned all the Elders, concrning leaving their calling to persue any occupation derogatory to that calling, assuring them that if persued, God would let them run themselves into difficulties, that he may stop them in their career, that salvation may come unto them, Elder Boyngton then arose and still attributed his difficulties \& conduct to the failure of the bank, stating that the bank he understood was instituted by the will \& revilations of God, \& he had been told that it never would fail let men do what they pleased,%
						\hypertarget{EMS-1837-JS-3-SEP-1837-b}%
					Pres.
						\marginnote{b}%
					Smith then arose, and stated that if this had been published, it was without authority, at least from him, he stated that he allways said that unless the institution was conducted upon righteous principles it could not stand, The church was then called upon to know whether they were sattisfied with the confession of Elder Boyngton, Voted in the negative Adjourned for one hour...
					
					...The Pres then arose \& made some remarks concerning the former presidents of the Seventies, their calling the authority of the priesthood \&.c...
					
		% -----
		\subsubsection{Letter to John Corrill and the Church in Missouri, 4 September 1837}
		% -----
			\label{EMS-1837-JS-4-SEP-1837}
			% \label{EMS-1830-JS-DAY-3LETTERMONTH-YEAR}
			
			\begin{description}
					\item[Print Sources]
				\end{description}

					\begin{tabular}{p{0.5in}p{1in}p{0.5in}p{1in}}
					JSP	 	& D5:426--431; J1:240--244	& JSR 		& --- \\
					DC	 	& ---				& HC 		& 2:508--509 \\
					HDDC 	& ---				& TCHC 		& 2:485 \\
					\end{tabular}
					% HDDC = woodford's DC diss; JSR = Marquardt's JS Revelations; TCHC = Vogel HC
				
				\begin{description}
					\item[Online Access] \href{https://www.josephsmithpapers.org/paper-summary/letter-to-john-corrill-and-the-church-in-missouri-4-september-1837/1}{Here}
					% \item[Analysis] \hyperref[XX]{Here}
				\end{description}
				
				\begin{description}
					\item[Background]
				\end{description}
				
					% It might also be useful to give a two sentence context of the period: e.g., "This speech was given in Nauvoo to a small congregation one month prior to the destruction of the Nauvoo Expositor. These and other tensions may have been on the prophet's mind when he gave the speech."
				
				\begin{description}
					\item[Original Source]
				\end{description}
				
					% brief description of original source material, with reference to JSP or somewhere else for more info
					% Background material: it might be helpful to have a note that says whether a given document was presented as a speech, was written in a journal, etc.
					% Also a comment here about how reliable the source might be, or what shape it is in, if there are large omissions or parts that are hard to understand, and so on.
				
				\begin{description}
					\item[Text]
				\end{description}
				
						\hypertarget{EMS-1837-JS-4-SEP-1837-a}%
					...Joseph
						\marginnote{a}%
					Smith J\supersu{r}\supers{.} Pres\supers{t} of the Church <of Christ> of Latter Day Saints in all the world
					
					To John Corroll [Corrill] \& the whole Church in Zion [p. 18] Sendeth greeting, Blessed be the God \sout{of} and father of our Lord Jesus Christ Who has blessed you with many blessings in Christ, And who has delivered you many times from the hands of your enimies And planted you many times in an heavenly or holy place, My respects \& love to you all, and my blessings upon all the faithfull \& true harted in the new \& everlasting covenant \& for as much as I have desired for a long time to see your faces, \& converse with you \& instruct you in those things which have been revealed to Me partaining to the Kingdom of God in the last days, I now write unto you offering an appolegy [apology], My being bound with bonds of affliction by the workers of iniquity and by the labours of the Church endeaveroung in all things to do the will of God, for the salvation of the Church both in temporal as well as spiritual things.
						\hypertarget{EMS-1837-JS-4-SEP-1837-b}%
					Bretheren
						\marginnote{b}%
					we have waided through a scene of affliction and sorrow thus far for the will of God, that language is inadequate to describe pray ye\sout{a} therefore with more earnestness for our redemption, You have undoubtedly been informed by letter \& otherwise of our difficulties in Kirtland which are now about being settled and that you may have a knowledge of the same I subscribe to you the following minuits of the comittee, of the whole Church of Kirtland the authorities \&.c. refering you to my brother Hyrum [Smith] \& br. T[homas] B. Marsh for further particulars also that you may know how to proceed to set in order \& regulate the affairs of the Church in zion whenever they become disorganized The minuits are as follows...
					
					...Pres\supersu{t}\supers{.} Smith then arose and stated that if this had been declared, no one had authority from him for so doing, For he had allways said unless the institution was conducted on richeous [righteous] principals it would not stand...
					
					...The Pres then arose and made some remarks concerning the formers Pres\supersu{ts} of the Seventies, the callings and authorities of their Priesthood \&c. \&c...
					
						\hypertarget{EMS-1837-JS-4-SEP-1837-c}%
					...Oliver
						\marginnote{c}%
					Cowdery has been in transgression, but as he is now chosen as one of the Presidents or councilors I trust that he will yet humble himself \& magnify his calling but if he should not, the church will soon be under the necessaty of raising their hands against him Therefore pray for him, David Whitmer Leonard Rich \& others have been in transgression but we hope that they may be humble \& ere long make sattisfaction to the Church otherwise they cannot retain their standing, Therefore we say unto you beware of all disaffected Characters for they come not to build up but to destroy \& scatter abroad, Though we or an Angel from Heaven preach any other Gospel or introduce [any other?] order of things <than> those things which ye have received and are authorized to receive\sout{d} from the first Presidency let him be accursed, May God Almighty Bless you all \& keep you unto the coming \& kingdom of our Lord and Savior Jesus Christ; Yours in the Bonds of the new <covenent>--- J. Smith, Jr....
					
		% -----
		\subsubsection{Revelation, 4 September 1837}
		% -----
			\label{EMS-1837-JS-4-SEP-1837-REV}
			% \label{EMS-1830-JS-DAY-3LETTERMONTH-YEAR}
			
			\begin{description}
					\item[Print Sources]
				\end{description}

					\begin{tabular}{p{0.5in}p{1in}p{0.5in}p{1in}}
					JSP	 	& D5:431--433; J1:245		& JSR 		& 283 \\
					DC	 	& ---				& HC 		& 2:511 \\
					HDDC 	& ---				& TCHC 		& 2:488 \\
					\end{tabular}
					% HDDC = woodford's DC diss; JSR = Marquardt's JS Revelations; TCHC = Vogel HC
				
				\begin{description}
					\item[Online Access] \href{https://www.josephsmithpapers.org/paper-summary/revelation-4-september-1837/1}{Here}
					% \item[Analysis] \hyperref[XX]{Here}
				\end{description}
				
				\begin{description}
					\item[Background]
				\end{description}
				
					% It might also be useful to give a two sentence context of the period: e.g., "This speech was given in Nauvoo to a small congregation one month prior to the destruction of the Nauvoo Expositor. These and other tensions may have been on the prophet's mind when he gave the speech."
				
				\begin{description}
					\item[Original Source]
				\end{description}
				
					% brief description of original source material, with reference to JSP or somewhere else for more info
					% Background material: it might be helpful to have a note that says whether a given document was presented as a speech, was written in a journal, etc.
					% Also a comment here about how reliable the source might be, or what shape it is in, if there are large omissions or parts that are hard to understand, and so on.
				
				\begin{description}
					\item[Text]
				\end{description}
				
					\begin{tightcenter}
					Revelation to Joseph Smith Jr Given in Kirtland Geauga Co. Ohio Sept 4\supers{th} 1837
					\end{tightcenter}

					Making known the transgression of John Whitmer W[illiam] W. Phelps

					Verily thus saith the Lord unto you my Servent Joseph. My Servents John Whitmer \& William W Phelps have done those things which are not pleasing in my sight Therefore if they repent not they shall be removed out of their places Amen---

					\begin{tightright}
					J Smith Jr
					\end{tightright}
					
		% -----
		\subsubsection{Minutes, 17 September 1837--A}
		% -----
			\label{EMS-1837-JS-17-SEP-1837-A}
			% \label{EMS-1830-JS-DAY-3LETTERMONTH-YEAR}
			
			\begin{description}
					\item[Print Sources]
				\end{description}

					\begin{tabular}{p{0.5in}p{1in}p{0.5in}p{1in}}
					JSP	 	& D5:442--443		& JSR 		& --- \\
					DC	 	& ---				& HC 		& --- \\
					HDDC 	& ---				& TCHC 		& --- \\
					\end{tabular}
					% HDDC = woodford's DC diss; JSR = Marquardt's JS Revelations; TCHC = Vogel HC
				
				\begin{description}
					\item[Online Access] \href{https://www.josephsmithpapers.org/paper-summary/minutes-17-september-1837-a/1}{Here}
					% \item[Analysis] \hyperref[XX]{Here}
				\end{description}
				
				\begin{description}
					\item[Background]
				\end{description}
				
					% It might also be useful to give a two sentence context of the period: e.g., "This speech was given in Nauvoo to a small congregation one month prior to the destruction of the Nauvoo Expositor. These and other tensions may have been on the prophet's mind when he gave the speech."
				
				\begin{description}
					\item[Original Source]
				\end{description}
				
					% brief description of original source material, with reference to JSP or somewhere else for more info
					% Background material: it might be helpful to have a note that says whether a given document was presented as a speech, was written in a journal, etc.
					% Also a comment here about how reliable the source might be, or what shape it is in, if there are large omissions or parts that are hard to understand, and so on.
				
				\begin{description}
					\item[Text]
				\end{description}
				
					...After some remarks by Presidents Smith \& Rigdon and others upon the disipline of Children, \&c. the administration of the Lords Supper being attended to, the meeting Closed by a benediction from the Bishop...
					
		% -----
		\subsubsection{Minutes, 17 September 1837--B}
		% -----
			\label{EMS-1837-JS-17-SEP-1837-B}
			% \label{EMS-1830-JS-DAY-3LETTERMONTH-YEAR}
			
			\begin{description}
					\item[Print Sources]
				\end{description}

					\begin{tabular}{p{0.5in}p{1in}p{0.5in}p{1in}}
					JSP	 	& D5:444--446		& JSR 		& --- \\
					DC	 	& ---				& HC 		& 2:513--514 \\
					HDDC 	& ---				& TCHC 		& 2:491--492 \\
					\end{tabular}
					% HDDC = woodford's DC diss; JSR = Marquardt's JS Revelations; TCHC = Vogel HC
				
				\begin{description}
					\item[Online Access] \href{https://www.josephsmithpapers.org/paper-summary/minutes-17-september-1837-b/1}{Here}
					% \item[Analysis] \hyperref[XX]{Here}
				\end{description}
				
				\begin{description}
					\item[Background]
				\end{description}
				
					% It might also be useful to give a two sentence context of the period: e.g., "This speech was given in Nauvoo to a small congregation one month prior to the destruction of the Nauvoo Expositor. These and other tensions may have been on the prophet's mind when he gave the speech."
				
				\begin{description}
					\item[Original Source]
				\end{description}
				
					% brief description of original source material, with reference to JSP or somewhere else for more info
					% Background material: it might be helpful to have a note that says whether a given document was presented as a speech, was written in a journal, etc.
					% Also a comment here about how reliable the source might be, or what shape it is in, if there are large omissions or parts that are hard to understand, and so on.
				
				\begin{description}
					\item[Text]
				\end{description}
				
					...Minuits of a conference of Elders held in the house of the Lord this evening Pres. Joseph Smith Jr presided, the conferance was op[e]ned by prayer by Pres. <S[idney]> Rigdon after which the conferance was addressed \sout{by} from the Chair, on the subject of the gathering of the Saints in the last days and the duties of the of the different quorums relations thereto, It appeared manifest to the conference that the places appointed for the gathering of the saints <were> at this time crowded to overflowing \& that it was necessary that there be more Stakes of Zion appointed in order that the poor might have a place to gather to, wherefore it was moved seconded \& carried by vote of the whole that Presidents J Smith Jr \& S. Rigdon be requested by this conference to go \& appoint other Stakes or places of gathering and that they receive a certificate of this their appointment signed by the Clerk of the Church...
					
		% -----
		\subsubsection{Travel Account and Questions, November 1837}
		% -----
			\label{EMS-1837-JS-NOV-1837}
			% \label{EMS-1830-JS-DAY-3LETTERMONTH-YEAR}
			
			\begin{description}
					\item[Print Sources]
				\end{description}

					\begin{tabular}{p{0.5in}p{1in}p{0.5in}p{1in}}
					JSP	 	& D5:478--484		& JSR 		& --- \\
					DC	 	& ---				& HC 		& --- \\
					HDDC 	& ---				& TCHC 		& --- \\
					\end{tabular}
					% HDDC = woodford's DC diss; JSR = Marquardt's JS Revelations; TCHC = Vogel HC
				
				\begin{description}
					\item[Online Access] \href{https://www.josephsmithpapers.org/paper-summary/travel-account-and-questions-november-1837/1}{Here}
					% \item[Analysis] \hyperref[XX]{Here}
				\end{description}
				
				\begin{description}
					\item[Background]
				\end{description}
				
					% It might also be useful to give a two sentence context of the period: e.g., "This speech was given in Nauvoo to a small congregation one month prior to the destruction of the Nauvoo Expositor. These and other tensions may have been on the prophet's mind when he gave the speech."
				
				\begin{description}
					\item[Original Source]
				\end{description}
				
					% brief description of original source material, with reference to JSP or somewhere else for more info
					% Background material: it might be helpful to have a note that says whether a given document was presented as a speech, was written in a journal, etc.
					% Also a comment here about how reliable the source might be, or what shape it is in, if there are large omissions or parts that are hard to understand, and so on.
				
				\begin{description}
					\item[Text]
				\end{description}
				
						\hypertarget{EMS-1837-JS-NOV-1837-a}%
					Be
						\marginnote{a}%
					it known unto the Saints scattered abroad greeting:

					That myself together with my beloved brother Sidney Rigdon, having been appointed by a general conference of elders held in Kirtland in the house of the Lord on the 18th of Sept. for the purpose of establishing places of gathering for the Saints \&c. we therefore would inform our readers that we started from Kirtland in company with V[inson] Knight and Wm. Smith on the 27th of Sept. last, for the purpose of vislting [visiting] the Far West, and also to discover situations suitable for the location of the Saints who are gathering for a refuge and safety, in the day of the wrath of God which is soon to burst upon the head of this generation, according to the testimony of the prophets; who speak expressly concerning the last days:
						\hypertarget{EMS-1837-JS-NOV-1837-b}%
					We
						\marginnote{b}%
					had a prosperous and a speedy journey; we held one meeting in Norton township Ohio, and three in Doublin, Ia. [Dublin, Indiana] one between Doublin and Ter[r]e Haute, Ia. two in Tere Haute, one in Palmyra, Mo. 2 in Huntsville, one in Carlton; all of which were tended with good success and generally allayed the prejudice and feeling of the people, as we judge from the treatment we received, being kindly and hospitably entertained. On our arrival at the city of Far West, we found the church of Latter Day Saints in that place in as prosperous a condition as we could have expected, and as we believe enjoying a goodly portion of the Spirit of God, to the joy and satisfaction of our hearts.

						\hypertarget{EMS-1837-JS-NOV-1837-c}%
					The
						\marginnote{c}%
					High council was immediately called and many difficulties adjusted, and the object of our mission was laid before them, after which the subject of the propriety of the Saints, gathering to the city of Far West, was taken into consider[a]tion, after a lengthy discussion upon the subject, it was voted, that the work of the gathering to that place be continued, and that there is a plenty of provisions in the upper counties for the support of that place, and also the emigration of the Saints; also voted that other Stakes be appointed in the regions round about, therefore a committee was appointed to locate the same; consisting of Oliver Cowdery, David Whitmer, John Corril[l], and Lyman Wight; who started on their mission before we left.

						\hypertarget{EMS-1837-JS-NOV-1837-d}%
					It
						\marginnote{d}%
					was also voted that the Saints be directed to those men for instruction concerning those places; and it may be expected that all the information necessary will be had from them concerning the location of those places, roads \&c. Now we would recommend to the Saints scattered abroad, that they make all possible exertions to gather themselves together unto those places; as peace, verily thus saith the Lord, peace shall soon be taken from the earth, and it has already began to be taken; for a lying spirit has gone out upon all the face of the earth and shall perplex the nations, and shall stur them up to anger against one another: for behold saith the Lord, very fierce and very terrible war is near at hand, even at your doors, therefore make haste saith the Lord O ye my people, and gather yourselres together and be at peace among yourselves, or there shall be no saf[e]ty for you.

						\hypertarget{EMS-1837-JS-NOV-1837-e}%
					We
						\marginnote{e}%
					furthermore say to those who wish to stop short of the city of Far West, to call on us and get information concerning those places of gathering: We would say to the Saints it is now a day of warning and not of many words; therefore, a word to the wise is sufficent. We exhort you to remember the words of the prophet Malichi which says, bring ye all the tithes into the store house that there may be meet in mine house, and prove me herewith saith the Lord of hosts, if I will not open you the windows of heaven and pour you out a blessing, that there shall not be room enough to receive it, and I will rebuke the devourer for your sake, and he shall not destroy the fruits of your ground, neither shall your vine cast her fruit before the time in the field, saith the Lord of hosts, and all nations shall call you blessed for ye shall be a delightsome land satth [saith] the Lord of hosts. We would also say to the Saints, that we were much pleased with the location of the Far West, and also the society of that place; and we purpose of locating our families in that place as soon as our circumstances will admit.

					We shall be under the necessity of observing brevity in this our communication for want of room to publish it, and we shall close after naming a few questions which are daily and hourly asked by all classes of people whilst we are traveling, and will answer them in our next.

						\hypertarget{EMS-1837-JS-NOV-1837-f}%
					Firstly,
						\marginnote{f}%
					Do you believe the bible?

					2nd, Wherein do you differ from other denominations?

					3rd, Will every body be damned but Mormons?

					4th, How and where did you obtain the book of Mormon?

					5th, Do you believe Jo Smith to be a prophet?

					6th, Do the Mormons believe in having all things common?

					7th, Do the Momons believe in having more wives than one?

					8th, Can they riase the dead?

					9th, What signs does Jo Smith give to establish his divine mission?

						\hypertarget{EMS-1837-JS-NOV-1837-g}%
					10th,
						\marginnote{g}%
					Was not Jo Smith a money digger?

					11th, Did he not Steal his wife?

					12th, Do the people have to give up their money when they join his church?

					13th, Are the Mormons Abolitionists?

					14th, Do they not stur up the Indians to war and to commit depredations?

					15th, Do the Mormons baptize in the name of Jo Smith?

					16th, If the Mormon doctrine is true, what have become of all that have died since the days of the apostles?

					17th Does not Jo Smith pretend to be Jesus Christ?

					18th, Is there any thing in the bible that liscences you to believe in revelation now days?

					19th, Is not the canon of the scriptures full?

					20th, What are the fundamental principles of your religion?

					The above questions are as many as we probably shall have room to answer in our next article, though many more may hereafter be asked and answered as circumstances will permit.
					
		% -----
		\subsubsection{Discourse, Early December 1837}
		% -----
			\label{EMS-1837-JS-DEC-1837}
			% \label{EMS-1830-JS-DAY-3LETTERMONTH-YEAR}
			
			\begin{description}
					\item[Print Sources]
				\end{description}

					\begin{tabular}{p{0.5in}p{1in}p{0.5in}p{1in}}
					JSP	 	& ---		& JSR 		& --- \\
					DC	 	& ---		& HC 		& --- \\
					HDDC 	& ---		& TCHC 		& --- \\
					\end{tabular}
					% HDDC = woodford's DC diss; JSR = Marquardt's JS Revelations; TCHC = Vogel HC
				
				\begin{description}
					\item[Online Access] \href{https://www.josephsmithpapers.org/paper-summary/discourse-early-december-1837/1}{Here}
					% \item[Analysis] \hyperref[XX]{Here}
				\end{description}
				
				\begin{description}
					\item[Background]
				\end{description}
				
					% It might also be useful to give a two sentence context of the period: e.g., "This speech was given in Nauvoo to a small congregation one month prior to the destruction of the Nauvoo Expositor. These and other tensions may have been on the prophet's mind when he gave the speech."
				
				\begin{description}
					\item[Original Source]
				\end{description}
				
					% brief description of original source material, with reference to JSP or somewhere else for more info
					% Background material: it might be helpful to have a note that says whether a given document was presented as a speech, was written in a journal, etc.
					% Also a comment here about how reliable the source might be, or what shape it is in, if there are large omissions or parts that are hard to understand, and so on.
				
				\begin{description}
					\item[Text]
				\end{description}
				
						\hypertarget{EMS-1837-JS-DEC-1837-a}%
					Joseph
						\marginnote{a}%
					then Spake as never man had Spoken before; Sho[w]ing the Love of God. Portraid the gethering to gether of Saints and Sin[n]ers. Drew a figer [figure] of the gospel net, in which the good would be with the bad; And of the tares gro[w]ing with the wheat. And of the Saints So living, that the teares seeing their good works; might be led to Glorify God. Joseph Said, ``The duty of the Saints is to bind--- to-gether Humanity. To throw round \uline{all} the \uline{arm} of \uline{Love}. Not to be to zelous; fore ``There is a way that seameth right unto man but the end \uline{thereof} is death''. Many take the \uline{influence} of \uline{zeal} to be the Spirit of God, but it is not. The evol Spirits often decieve; and lead man astray. They are trying to divide the Saints.
					
						\hypertarget{EMS-1837-JS-DEC-1837-b}%
					But
						\marginnote{b}%
					we are cal[le]d-upon to hunt for the one that is mising from the fold, whether Saint or Siner. And now Bretheren I wish to Say something about the Bruster [James Brewster] and Noris [Moses Norris] Brothers and Sisters. I met with them before going to Missouri. The Spirit of \uline{Love}, Joy and peace rested on us, and the \uline{Holy} \uline{Ghost} bore witness of their faith in God; and \uline{now} I \uline{learn} they are all Cut off from the Church.
					
					I want some one to make a motion that all of them be restord to full fellowship (with out baptism) as though no action had ben taken against them''.

% -----------------------------------------------------------
\pagebreak{}
% -----------------------------------------------------------